\documentclass[12pt,a4paper]{article}
\usepackage[utf8]{inputenc}
\usepackage{amsmath}
\usepackage{amssymb}
\usepackage{amsthm}
\usepackage{algorithm}
\usepackage{algpseudocode}
\usepackage{graphicx}
\usepackage{hyperref}
\usepackage{geometry}
\geometry{margin=1in}

\title{The 8-Puzzle Problem: \\
       Hill Climbing and A* Search Algorithms}
\author{AI Lab Assignment}
\date{\today}

\newtheorem{definition}{Definition}
\newtheorem{theorem}{Theorem}
\newtheorem{lemma}{Lemma}

\begin{document}

\maketitle

\begin{abstract}
This document provides a comprehensive mathematical treatment of the 8-puzzle problem and its solution using Hill Climbing and A* search algorithms. We present formal definitions, algorithm pseudocode, heuristic functions, complexity analysis, and experimental results.
\end{abstract}

\tableofcontents
\newpage

\section{Introduction}

The 8-puzzle problem is a classic problem in artificial intelligence that involves sliding numbered tiles on a $3 \times 3$ grid to reach a goal configuration. This problem serves as an excellent testbed for studying search algorithms and heuristic functions.

\subsection{Problem Statement}

\begin{definition}[8-Puzzle State]
An 8-puzzle state $S$ is a $3 \times 3$ matrix where each cell contains either a number from $\{1, 2, 3, 4, 5, 6, 7, 8\}$ or a blank space (denoted as 0):
\[
S = \begin{bmatrix}
s_{11} & s_{12} & s_{13} \\
s_{21} & s_{22} & s_{23} \\
s_{31} & s_{32} & s_{33}
\end{bmatrix}
\]
where $s_{ij} \in \{0, 1, 2, 3, 4, 5, 6, 7, 8\}$ and each number appears exactly once.
\end{definition}

\begin{definition}[Valid Move]
A valid move $m$ slides a tile adjacent to the blank space into the blank space. The set of valid moves from state $S$ is:
\[
M(S) = \{m : m \in \{\text{Up, Down, Left, Right}\} \land \text{move } m \text{ is within bounds}\}
\]
\end{definition}

\begin{definition}[State Space]
The state space $\mathcal{S}$ is the set of all reachable states from an initial state $S_0$ through valid moves. The state space forms a directed graph $G = (\mathcal{S}, E)$ where:
\begin{itemize}
    \item $\mathcal{S}$ is the set of all possible board configurations
    \item $E = \{(S_i, S_j) : S_j$ can be reached from $S_i$ by one valid move$\}$
\end{itemize}
\end{definition}

\section{Heuristic Functions}

Heuristic functions estimate the cost from a current state to the goal state. A good heuristic function is both \textit{admissible} (never overestimates) and \textit{informative} (provides useful guidance).

\subsection{Manhattan Distance Heuristic}

\begin{definition}[Manhattan Distance]
For a state $S$, the Manhattan distance heuristic $h_M(S)$ is the sum of the distances of each tile from its goal position:
\[
h_M(S) = \sum_{i=1}^{8} |x_i - x_i^*| + |y_i - y_i^*|
\]
where $(x_i, y_i)$ is the current position of tile $i$ and $(x_i^*, y_i^*)$ is its goal position.
\end{definition}

\begin{theorem}[Admissibility of Manhattan Distance]
The Manhattan distance heuristic $h_M$ is admissible for the 8-puzzle problem.
\end{theorem}

\begin{proof}
Each tile must move at least the Manhattan distance to reach its goal position. Since each move changes the position of exactly one tile by one unit, the Manhattan distance never overestimates the minimum number of moves required.
\end{proof}

\subsection{Misplaced Tiles Heuristic}

\begin{definition}[Misplaced Tiles]
The misplaced tiles heuristic $h_T(S)$ counts the number of tiles not in their goal position:
\[
h_T(S) = |\{i : s_i \neq s_i^* \land s_i \neq 0\}|
\]
where $s_i$ is the value at position $i$ in state $S$ and $s_i^*$ is the goal value.
\end{definition}

\begin{theorem}[Admissibility of Misplaced Tiles]
The misplaced tiles heuristic $h_T$ is admissible for the 8-puzzle problem.
\end{theorem}

\begin{proof}
Each misplaced tile requires at least one move to reach its goal position. Therefore, $h_T(S)$ never overestimates the actual cost.
\end{proof}

\subsection{Dominance Relation}

\begin{theorem}[Heuristic Dominance]
For all non-goal states $S$ in the 8-puzzle:
\[
h_M(S) \geq h_T(S)
\]
Thus, Manhattan distance dominates the misplaced tiles heuristic.
\end{theorem}

\begin{proof}
For any misplaced tile, its Manhattan distance from the goal is at least 1. Therefore:
\[
h_M(S) = \sum_{i \in \text{misplaced}} d_M(i) \geq \sum_{i \in \text{misplaced}} 1 = h_T(S)
\]
\end{proof}

\section{Hill Climbing Algorithm}

Hill Climbing is a local search algorithm that iteratively moves to neighboring states with better heuristic values.

\subsection{Basic Hill Climbing}

\begin{algorithm}
\caption{Hill Climbing Search}
\begin{algorithmic}[1]
\Procedure{HillClimbing}{$S_0, S_{goal}, h$}
    \State $current \gets S_0$
    \State $visited \gets \{S_0\}$
    \While{$current \neq S_{goal}$}
        \State $neighbors \gets \text{GetSuccessors}(current) \setminus visited$
        \If{$neighbors = \emptyset$}
            \State \Return \textsc{Failure} \Comment{Stuck at local optimum}
        \EndIf
        \State $best \gets \arg\min_{n \in neighbors} h(n)$
        \If{$h(best) \geq h(current)$}
            \State \Return \textsc{Failure} \Comment{No improvement}
        \EndIf
        \State $current \gets best$
        \State $visited \gets visited \cup \{best\}$
    \EndWhile
    \State \Return $current$
\EndProcedure
\end{algorithmic}
\end{algorithm}

\subsection{Hill Climbing with Sideways Moves}

To escape plateaus, we allow a limited number of sideways moves (moves to states with equal heuristic value).

\begin{algorithm}
\caption{Hill Climbing with Sideways Moves}
\begin{algorithmic}[1]
\Procedure{HillClimbingSideways}{$S_0, S_{goal}, h, max_{sideways}$}
    \State $current \gets S_0$
    \State $visited \gets \{S_0\}$
    \State $sideways\_count \gets 0$
    \While{$current \neq S_{goal}$}
        \State $neighbors \gets \text{GetSuccessors}(current) \setminus visited$
        \If{$neighbors = \emptyset$}
            \State \Return \textsc{Failure}
        \EndIf
        \State $best \gets \arg\min_{n \in neighbors} h(n)$
        \If{$h(best) < h(current)$}
            \State $current \gets best$
            \State $sideways\_count \gets 0$
        \ElsIf{$h(best) = h(current)$ \textbf{and} $sideways\_count < max_{sideways}$}
            \State $current \gets best$
            \State $sideways\_count \gets sideways\_count + 1$
        \Else
            \State \Return \textsc{Failure}
        \EndIf
        \State $visited \gets visited \cup \{best\}$
    \EndWhile
    \State \Return $current$
\EndProcedure
\end{algorithmic}
\end{algorithm}

\subsection{Complexity Analysis}

\begin{theorem}[Time Complexity of Hill Climbing]
In the worst case, Hill Climbing explores $O(b^d)$ nodes, where $b$ is the branching factor and $d$ is the solution depth. However, in practice, it typically explores $O(d)$ nodes for problems without local optima.
\end{theorem}

\begin{theorem}[Space Complexity]
Hill Climbing has space complexity $O(|visited|) = O(d)$ in the best case, where $d$ is the number of steps taken.
\end{theorem}

\subsection{Limitations}

Hill Climbing suffers from three main problems:
\begin{enumerate}
    \item \textbf{Local Optima}: Getting stuck at states where all neighbors have worse heuristic values
    \item \textbf{Plateaus}: Flat regions where neighbors have equal heuristic values
    \item \textbf{Ridges}: Sequences of local maxima difficult to navigate
\end{enumerate}

\section{A* Search Algorithm}

A* is an informed search algorithm that combines actual cost and heuristic estimates to find optimal solutions.

\subsection{Cost Function}

\begin{definition}[A* Evaluation Function]
For a state $S$, A* uses the evaluation function:
\[
f(S) = g(S) + h(S)
\]
where:
\begin{itemize}
    \item $g(S)$ is the actual cost from the start state to $S$
    \item $h(S)$ is the heuristic estimate from $S$ to the goal
    \item $f(S)$ is the estimated total cost through $S$
\end{itemize}
\end{definition}

\subsection{Algorithm}

\begin{algorithm}
\caption{A* Search}
\begin{algorithmic}[1]
\Procedure{AStar}{$S_0, S_{goal}, h$}
    \State $openSet \gets$ priority queue containing $S_0$
    \State $closedSet \gets \emptyset$
    \State $g(S_0) \gets 0$
    \State $f(S_0) \gets h(S_0)$
    \While{$openSet \neq \emptyset$}
        \State $current \gets$ node in $openSet$ with minimum $f$ value
        \If{$current = S_{goal}$}
            \State \Return \textsc{ReconstructPath}($current$)
        \EndIf
        \State Remove $current$ from $openSet$
        \State Add $current$ to $closedSet$
        \For{each $neighbor$ in GetSuccessors($current$)}
            \If{$neighbor \in closedSet$}
                \State \textbf{continue}
            \EndIf
            \State $tentative\_g \gets g(current) + 1$
            \If{$neighbor \notin openSet$ \textbf{or} $tentative\_g < g(neighbor)$}
                \State $g(neighbor) \gets tentative\_g$
                \State $f(neighbor) \gets g(neighbor) + h(neighbor)$
                \State $parent(neighbor) \gets current$
                \If{$neighbor \notin openSet$}
                    \State Add $neighbor$ to $openSet$
                \EndIf
            \EndIf
        \EndFor
    \EndWhile
    \State \Return \textsc{Failure}
\EndProcedure
\end{algorithmic}
\end{algorithm}

\subsection{Optimality and Completeness}

\begin{theorem}[Optimality of A*]
If the heuristic function $h$ is admissible, A* is guaranteed to find an optimal solution.
\end{theorem}

\begin{theorem}[Completeness of A*]
A* is complete: if a solution exists, A* will find it.
\end{theorem}

\subsection{Complexity Analysis}

\begin{theorem}[Time Complexity of A*]
The time complexity of A* is $O(b^d)$ in the worst case, where $b$ is the branching factor and $d$ is the solution depth. With a good heuristic, this can be reduced significantly.
\end{theorem}

\begin{theorem}[Space Complexity of A*]
A* has space complexity $O(b^d)$ as it stores all generated nodes.
\end{theorem}

\section{Experimental Results}

\subsection{Test Cases}

We evaluated both algorithms on several 8-puzzle instances:

\subsubsection{Test Case 1: Simple Problem}

\textbf{Initial State:}
\[
S_0 = \begin{bmatrix}
1 & 2 & 3 \\
8 & 4 & 0 \\
7 & 6 & 5
\end{bmatrix}
\]

\textbf{Goal State:}
\[
S_{goal} = \begin{bmatrix}
1 & 2 & 3 \\
8 & 0 & 4 \\
7 & 6 & 5
\end{bmatrix}
\]

\textbf{Results:}
\begin{itemize}
    \item Hill Climbing (Manhattan): 1 move, 1 node expanded
    \item A* Search (Manhattan): 1 move, 1 node expanded
\end{itemize}

\subsubsection{Test Case 2: Medium Complexity}

\textbf{Initial State:}
\[
S_0 = \begin{bmatrix}
1 & 2 & 3 \\
0 & 8 & 4 \\
7 & 6 & 5
\end{bmatrix}
\]

\textbf{Goal State:}
\[
S_{goal} = \begin{bmatrix}
1 & 2 & 3 \\
8 & 0 & 4 \\
7 & 6 & 5
\end{bmatrix}
\]

\textbf{Results:}
\begin{itemize}
    \item Hill Climbing (Manhattan): 1 move, 1 node expanded
    \item A* Search (Manhattan): 1 move, 1 node expanded
\end{itemize}

\subsection{Performance Comparison}

\begin{table}[h]
\centering
\begin{tabular}{|l|c|c|}
\hline
\textbf{Algorithm} & \textbf{Moves} & \textbf{Nodes Expanded} \\
\hline
Hill Climbing (Manhattan) & 1 & 1 \\
Hill Climbing (Misplaced) & Variable & Variable \\
A* (Manhattan) & 1 & 1 \\
A* (Misplaced) & 1 & Variable \\
\hline
\end{tabular}
\caption{Algorithm comparison on simple test cases}
\end{table}

\subsection{Analysis}

\textbf{Advantages of Hill Climbing:}
\begin{itemize}
    \item Low memory requirements: $O(d)$
    \item Fast execution for problems without local optima
    \item Simple implementation
\end{itemize}

\textbf{Disadvantages of Hill Climbing:}
\begin{itemize}
    \item Can get stuck at local optima
    \item Not guaranteed to find a solution
    \item Not guaranteed to find optimal solution
\end{itemize}

\textbf{Advantages of A*:}
\begin{itemize}
    \item Guaranteed to find optimal solution (with admissible heuristic)
    \item Complete: always finds solution if one exists
    \item Efficient with good heuristics
\end{itemize}

\textbf{Disadvantages of A*:}
\begin{itemize}
    \item High memory requirements: $O(b^d)$
    \item Can be slow for complex problems
\end{itemize}

\section{Mathematical Properties}

\subsection{State Space Size}

\begin{theorem}[8-Puzzle State Space]
The total number of possible configurations in the 8-puzzle is $9! = 362,880$. However, only half are reachable from any given initial state due to parity constraints.
\end{theorem}

\begin{proof}
There are $9!$ ways to arrange 9 tiles (including the blank). The 8-puzzle can be proven to have two disjoint classes of states based on permutation parity, making only $9!/2 = 181,440$ states reachable from any given initial state.
\end{proof}

\subsection{Branching Factor}

\begin{lemma}[Average Branching Factor]
The average branching factor of the 8-puzzle is approximately 2.67.
\end{lemma}

\begin{proof}
The blank can be in one of 9 positions:
\begin{itemize}
    \item Corner positions (4): 2 possible moves each
    \item Edge positions (4): 3 possible moves each
    \item Center position (1): 4 possible moves
\end{itemize}
Average: $\frac{4(2) + 4(3) + 1(4)}{9} = \frac{24}{9} \approx 2.67$
\end{proof}

\section{Conclusion}

We have presented a comprehensive mathematical treatment of the 8-puzzle problem and two solution approaches: Hill Climbing and A* search. While Hill Climbing offers simplicity and efficiency for easy problems, A* provides optimality guarantees at the cost of higher memory usage.

The choice between algorithms depends on the problem characteristics:
\begin{itemize}
    \item Use Hill Climbing for: simple problems, memory-constrained environments, quick approximate solutions
    \item Use A* for: problems requiring optimal solutions, when memory is available, guaranteed completeness
\end{itemize}

\section{References}

\begin{enumerate}
    \item Russell, S., \& Norvig, P. (2020). \textit{Artificial Intelligence: A Modern Approach} (4th ed.). Pearson.
    \item Pearl, J. (1984). \textit{Heuristics: Intelligent Search Strategies for Computer Problem Solving}. Addison-Wesley.
    \item Korf, R. E. (1985). Depth-first iterative-deepening: An optimal admissible tree search. \textit{Artificial Intelligence}, 27(1), 97-109.
\end{enumerate}

\end{document}
